% TOP_PROBLEM: {"id": 30006957, "title": "David–Semmes Riesz-transform problem in intermediate dimensions", "category_id": 8, "category": {"id": 8, "name": "analysis", "display_name": "Analysis", "description": "Limits, continuity, calculus, and function theory.", "slug": "analysis", "order_index": 8, "created_at": "2026-07-31T15:26:25.671Z"}, "status": "open"}
% FIRST_POSTED: 2026-09-25
% ENTRY_KIND: statement_only
% !TeX program = lualatex
% Definition and sources only; this file is not an LLM solution attempt.
\documentclass[11pt]{article}
\usepackage[margin=25mm]{geometry}
\usepackage{fontspec}
\setmainfont{Latin Modern Roman}
\usepackage{amsmath,amssymb,mathtools}
\usepackage{unicode-math}
\setmathfont{Latin Modern Math}
\usepackage{xurl,hyperref}
\hypersetup{colorlinks=true,urlcolor=blue}
\setcounter{secnumdepth}{0}
\setlength{\parindent}{0pt}
\setlength{\parskip}{5pt}
\setlength{\emergencystretch}{3em}
\begin{document}
\section*{\texorpdfstring{David–Semmes Riesz-transform problem in intermediate dimensions}{David–Semmes Riesz-transform problem in intermediate dimensions}}\label{30006957}
\subsection{Definitions and mathematical statement}
A Radon measure $\mu$ on ${\mathbb{R}}^d$ is $n$-Ahlfors--David regular if some $C\ge1$ satisfies
\[
 C^{-1}r^n\le\mu(B(x,r))\le Cr^n
\]
for $x\in\operatorname{supp}\mu$ and $0<r<\operatorname{diam}(\operatorname{supp}\mu)$. Uniform $n$-rectifiability can be defined by big pieces of Lipschitz images: fixed $\theta,L>0$ work in every such ball, giving an $L$-Lipschitz map from an $n$-ball of radius $r$ whose image meets that ball in $\mu$-measure at least $\theta r^n$.

For $2\le n\le d-2$, suppose the truncated vector Riesz transforms
\[
 R_{\mu,{\varepsilon}}f(x)=\int_{|x-y|>{\varepsilon}}
 \frac{x-y}{|x-y|^{n+1}}f(y){\,\mathrm{d}}\mu(y)
\]
have uniformly bounded norms from $L^2(\mu)$ to $L^2(\mu;{\mathbb{R}}^d)$ over all ${\varepsilon}>0$. Must $\mu$ be uniformly $n$-rectifiable? The constants need not be numerically prescribed; they must be uniform across location and scale.
\par\smallskip\textit{Formulation and concept references:} \href{https://arxiv.org/html/2509.26547}{[S1]}.

\subsection{Short English statement}
Does square-integrable boundedness of the Riesz transform force quantitative rectifiability for regular measures in the intermediate dimensions?
\subsection{Sources}
\begin{itemize}
\item[S1]Xavier Tolsa, Riesz transforms and the BAUPP and BWGL criteria for uniform rectifiability (2025).
\url{https://arxiv.org/html/2509.26547}.
\textit{Formulation source.}
\end{itemize}
\end{document}
